\chapter{الطرح}\label{ch:g2:subtraction}

حين تكون آحاد العدد الأعلى أقلّ من أن نطرح منها، نفكّ عشرة ---
وهذا هو \emph{الاستلاف}، مرآة الاحتفاظ. يشرحه هذا الفصل
بالأعواد، ثم يدرّب على الطرح المكتوب، مع العدّ إلى الأمام
مختصرًا لطيفًا.

\section{لماذا نستلف}

\begin{example}[فكّ عشرة]\label{ex:g2:subtraction:unbundle}
$42 - 17$: نريد أن نأخذ $7$ آحاد من $2$ من الآحاد --- وهي أقلّ
من ذلك! فنفكّ واحدة من العشرات $4$ إلى عشرة آحاد مفردة: فيصير
الأعلى $3$ عشرات و $12$ واحدًا. ثم نأخذ: $12 - 7 = 5$ آحاد، و
$3 - 1 = 2$ من العشرات. إذن $42 - 17 = 25$.
\end{example}

\begin{omfigure}
\begin{tikzpicture}[scale=0.75]
  % عشرة مربوطة في حزمة
  \foreach \i in {0,...,9} \draw[omExo, very thick]
    (0.2*\i,0) -- (0.2*\i,1.7);
  \draw[omExo, thick, rounded corners=3pt]
    (-0.2,0.62) rectangle (2.0,1.05);
  \node[below, font=\small] at (0.9,-0.25) {عشرة واحدة};
  % ثم تُفكّ إلى عشرة آحاد مفردة
  \draw[-Stealth, omThm, very thick] (2.6,0.85) -- (3.6,0.85);
  \foreach \i in {0,...,9} \draw[omDef, very thick]
    (4.2+0.34*\i,0) -- (4.2+0.34*\i,1.7);
  \node[below, font=\small] at (5.73,-0.25) {عشرة آحاد مفردة};
  \node[right, font=\small, align=left] at (8.0,0.85)
    {فُكّ عشرة ليكفي عدد\\الآحاد للطرح};
\end{tikzpicture}

{\small الاستلاف فكّ للحزمة: عشرة واحدة تصير عشرة آحاد، فيمضي
الطرح.}
\end{omfigure}

\section{الطريقة المكتوبة}

\begin{method}[الطرح بالأعمدة]\label{met:g2:subtraction:column}
\begin{enumerate}
  \item اُكتب العدد الأكبر في الأعلى، محاذيًا من اليمين؛
  \item اِطرح الآحاد؛ فإن كان الرقم الأعلى أصغر من أن يكفي،
        فاستلف عشرة (فتصير عشرة آحاد في العمود الحالي)؛
  \item اِطرح العشرات، ثم المئات؛
  \item \emph{تحقّق}: الناتج زائد المطروح يجب أن يعيد العدد
        الأعلى.
\end{enumerate}
\end{method}

\begin{example}\label{ex:g2:subtraction:worked}
احسب $503 - 168$:
\[
\begin{array}{r}
5\,0\,3 \\
-\ 1\,6\,8 \\
\hline
3\,3\,5
\end{array}
\]
الآحاد: $3 - 8$ غير ممكن؛ فاستلف عبر العشرات (ولا عشرات هنا،
فاستلف من المئات أولًا): فيصير الأعلى $4$ مئات و
$9$ عشرات و $13$ واحدًا. ثم $13 - 8 = 5$؛ و $9 - 6 = 3$؛
و $4 - 1 = 3$. والتحقّق: $335 + 168 = 503$. \checkmark
\end{example}

\section{العدّ إلى الأمام}

\begin{method}[الطرح بالعدّ إلى الأمام]\label{met:g2:subtraction:countup}
في الأعداد المتقاربة، اِصعد من الصغير إلى الكبير واجمع
القفزات:
\[
83 - 78: \quad 78 \to 80 \text{ يساوي } 2, \ 80 \to 83 \text{ يساوي } 3,
\text{ إذن } 83 - 78 = 5 .
\]
والعدّ إلى الأمام أسرع من الاستلاف حين يتقارب العددان --- وهو
تمامًا ما يفعله البائع حين يردّ الباقي.
\end{method}

\begin{omfigure}
\begin{tikzpicture}[scale=0.85]
  \draw[-Stealth] (76.6,0) -- (84.4,0);
  \foreach \i in {77,...,84} \draw (\i,-0.1) -- (\i,0.1);
  \foreach \i in {78,80,83}
    \node[below, font=\footnotesize] at (\i,-0.12) {\i};
  \draw[-Stealth, omDef, thick] (78,0.3) .. controls (79,0.9) ..
    (80,0.3);
  \node[omDef, font=\small] at (79,1.0) {$+2$};
  \draw[-Stealth, omThm, thick] (80,0.3) .. controls (81.5,1.05) ..
    (83,0.3);
  \node[omThm, font=\small] at (81.5,1.1) {$+3$};
\end{tikzpicture}

{\small $83 - 78$ بالعدّ إلى الأمام: قفزتان، $2 + 3 = 5$.}
\end{omfigure}

\begin{example}[ردّ الباقي]\label{ex:g2:subtraction:change}
تكلّف لعبة $34$؛ وتدفع بورقة $50$. والباقي: عُدّ إلى الأمام،
$34 \to 40$ مقداره $6$، و $40 \to 50$ مقداره $10$، فالباقي $16$.
وهو الطرح $50 - 34 = 16$، بطريقة البائع.
\end{example}

\section{تمارين}

\begin{exercise}[$\star$]\label{exo:g2:subtraction:1}
احسب (بلا استلاف): $58 - 23$؛ \quad $176 - 45$؛ \quad
$389 - 164$.
\end{exercise}

\begin{exercise}[$\star$]\label{exo:g2:subtraction:2}
احسب مع الاستلاف: $52 - 17$؛ \quad $73 - 28$؛ \quad
$60 - 24$.
\end{exercise}

\begin{exercise}[$\star$]\label{exo:g2:subtraction:3}
احسب بالأعمدة وتحقّق من كل عملية: $324 - 167$؛ \quad
$500 - 236$؛ \quad $703 - 158$.
\end{exercise}

\begin{exercise}[$\star$]\label{exo:g2:subtraction:4}
احسب بالعدّ إلى الأمام: $62 - 58$؛ \quad $91 - 85$؛ \quad
$100 - 96$.
\end{exercise}

\begin{exercise}[$\star$]\label{exo:g2:subtraction:5}
احسب الباقي (بالعدّ إلى الأمام): ادفع $50$ ثمنًا لكتاب قيمته $42$؛
وادفع $20$ ثمنًا للعبة قيمتها $13$؛ وادفع $100$ ثمنًا لحقيبة قيمتها $75$.
\end{exercise}

\begin{exercise}[$\star$]\label{exo:g2:subtraction:6}
«طول شجرة $156$ سم. وقد نمت $28$ سم هذه السنة. فكم كان طولها
السنة الماضية؟» اُكتب عملية الطرح والجواب.
\end{exercise}

\begin{exercise}[$\star$]\label{exo:g2:subtraction:7}
تحقّق بالجمع التوأم: $85 - 39 = 46$؛ أصحيح هذا؟
و $214 - 68 = 156$؛ أصحيح هذا؟
\end{exercise}

\begin{exercise}[$\star$]\label{exo:g2:subtraction:8}
لكل قصة، اختر $+$ أو $-$: «$47$ طائرًا، طار منها $19$»؛
«$47$ تفاحة، وقُطف $19$ غيرها»؛ «عند سامي $47$، وعند لينا $19$:
بكم يزيد سامي؟».
\end{exercise}

\begin{exercise}[$\star$]\label{exo:g2:subtraction:9}
أكمل الفراغ: $63 - \;?\; = 45$. (عُدّ إلى الأمام من $45$ إلى $63$.)
\end{exercise}

\begin{exercise}[$\star\star$]\label{exo:g2:subtraction:10}
لعبة: فكّر في عدد، ثم أضف إليه $30$، ثم اطرح منه $18$.
فتكون النتيجة $58$. فما العدد الأول؟ (اِعمل بالعكس!)
\end{exercise}
